\documentclass[journal]{IEEEtran}
\usepackage{amsmath,amssymb,bm}
\usepackage{amsthm}
\usepackage{graphicx}
\usepackage[caption=false,font=footnotesize]{subfig}
\usepackage{booktabs}
\usepackage{cite}
\usepackage[hidelinks]{hyperref}
\usepackage{url}
\theoremstyle{definition}
\newtheorem{assumption}{Assumption}[section]
\newtheorem{definition}{Definition}[section]
\newtheorem{problem}{Problem}[section]
\theoremstyle{plain}
\newtheorem{theorem}{Theorem}[section]
\newtheorem{lemma}{Lemma}[section]
\newtheorem{proposition}{Proposition}[section]
\newtheorem{corollary}{Corollary}[section]
\theoremstyle{remark}\newtheorem{remark}{Remark}[section]


\title{Prescribed-Time Trajectory Tracking Control of a Free-Floating Space Manipulator Based on Smooth Periodic Delayed Feedback}
\author{Author~Name\\
\thanks{The authors are with the Department of XXX, XXX University, XXX (e-mail: author@xxx.edu).}}


\begin{document}\maketitle
\begin{abstract}
This paper proposes a joint-space trajectory-tracking method based on smooth periodic delayed feedback (PDF) for a pre-capture free-floating space manipulator whose base is uncontrolled, whose dynamics are strongly coupled, and whose end-effector tracking accuracy is demanding. First, momentum conservation is used to reduce the dynamics with base reaction motion to an equivalent joint-space model. For bounded lumped disturbances, a prescribed-time disturbance observer generates $\hat d(t)$ and compensates matched disturbances after a specified observation time. A dynamic compensation term then establishes an auxiliary input channel with a controllable canonical form. Finally, periodic time-varying delayed feedback yields a prescribed-time tracking law using both current and delayed errors. Under continuous-time implementation, an exact model, and exact compensation, the closed-loop error becomes exactly zero after the prescribed time $T_o+2h$. With observation error, model approximation, actuator limits, and numerical discretization, this result becomes practical bounded tracking near the prescribed time, with explicit error bounds. MATLAB simulations using a nominal equivalent joint-space model parameterized for a seven-degree-of-freedom space manipulator demonstrate numerical feasibility under bounded disturbances.
\end{abstract}
\begin{IEEEkeywords}
Free-floating space manipulator, trajectory tracking, prescribed-time control, periodic delayed feedback, disturbance observer.
\end{IEEEkeywords}


\section{Introduction}
Space manipulators are central to on-orbit servicing, debris removal, and maintenance of failed spacecraft. Unlike ground-based manipulators, a pre-capture space manipulator is free floating: joint motion induces base reaction motion through momentum conservation and changes the end-effector trajectory. Its dynamics are nonlinear and strongly coupled, with parameter uncertainty, time-varying disturbances, and actuator faults. Accurate tracking with a prescribed convergence time is therefore a fundamental problem.

\subsection{Related Work and Background}
Under zero initial linear and angular momentum, base velocity is determined by joint velocity, and end-effector velocity is described by the generalized Jacobian rather than a fixed-base Jacobian. The generalized Jacobian was introduced by Umetani and Yoshida~\cite{umetani1989}; the virtual manipulator approach characterized the coupled kinematics and dynamics~\cite{vafa1990}. Dynamic singularities and the properties of free-floating and free-flying robots were studied in~\cite{papadopoulos1993,dubowsky1993}. Adaptive computed-torque and reactionless-motion controllers were developed in~\cite{parlaktuna2004,xu2016}. In particular, the reduced dynamics can be written as
\[
H_g(q)\ddot q+C_g(q,\dot q)\dot q=\tau_m.
\]
Trajectory planning, model predictive control, time-delay estimation, adaptive sliding-mode control, and prescribed-performance control have also been investigated~\cite{wang2015,ning2019,yu2022,yao2021}. Their convergence times, however, generally depend on gains and initial conditions.

Fixed-time and prescribed-time methods provide convergence bounds selected at design time~\cite{polyakov2012,zuo2019}. Robot tracking and space-robot pose control have been addressed using these ideas~\cite{su2020,xuPredefinedTimeTrackingControl2024,li2025}, but nonsmooth terminal surfaces or fractional powers can cause input jumps or gain singularities. Pyragas introduced delayed feedback control~\cite{pyragas1992}; smooth periodic delayed feedback (PDF) provides fixed-time stabilization for controllable linear time-delay systems while avoiding terminal gain singularities~\cite{zhouFixedtimeStabilizationLinear2022}. Its mechanism is the annihilation or nilpotence of a finite-dimensional one-period state-transition matrix. Its use for nonlinear manipulators requires accurate disturbance compensation. Prescribed-time disturbance observers provide an appropriate interface~\cite{jiangPrescribedtimeDisturbanceObserver2024}.

\subsection{Contributions and Organization}
This paper develops a four-stage structure: disturbance observation, dynamic compensation, error linearization, and periodic delayed feedback. The contributions are: (1) a formal reduction from base reaction motion and momentum conservation to equivalent joint-space dynamics; (2) a prescribed-time disturbance observer and a two-stage timing design; (3) a normalized controllable error channel for PDF design; (4) a Gramian-based gain construction and explicit practical error bounds; and (5) MATLAB validation with a seven-degree-of-freedom model.

The remainder is organized as follows. Section~\ref{sec:modeling} establishes the model; Section~\ref{sec:problem} formulates the problem; Section~\ref{sec:controller} designs the observer, compensation, and PDF controller; Section~\ref{sec:stability} proves convergence; Section~\ref{sec:simulation} presents simulations; and Section~\ref{sec:conclusion} concludes the paper.

\section{Modeling of a Free-Floating Space Manipulator}
\label{sec:modeling}
This section establishes the equivalent joint-space dynamics of a pre-capture free-floating space manipulator. Momentum conservation eliminates the base acceleration and yields a reduced model suitable for controller design.
\subsection{Notation}
Consider a free-floating manipulator with a base and $n$ revolute joints. Let $\dot x_b=[v_b^\top,\ \omega_b^\top]^\top\in\mathbb R^6$, $q=[q_1,\ldots,q_n]^\top\in\mathbb R^n$, and $\dot x_e=[v_e^\top,\ \omega_e^\top]^\top\in\mathbb R^6$.
\subsection{Full Dynamics}
With no external force or torque, the kinetic energy is
\begin{equation}T=\frac12\begin{bmatrix}\dot x_b^\top & \dot q^\top\end{bmatrix}\begin{bmatrix}H_b(q)&H_{bm}(q)\\H_{bm}^\top(q)&H_m(q)\end{bmatrix}\begin{bmatrix}\dot x_b\\\dot q\end{bmatrix}.\label{eq:kinetic_energy_ffsm}\end{equation}
The Lagrange equations are
\begin{equation}\begin{bmatrix}H_b&H_{bm}\\H_{bm}^\top&H_m\end{bmatrix}\begin{bmatrix}\ddot x_b\\\ddot q\end{bmatrix}+\begin{bmatrix}C_b&C_{bm}\\C_{bm}^\top&C_m\end{bmatrix}\begin{bmatrix}\dot x_b\\\dot q\end{bmatrix}=\begin{bmatrix}0\\\tau\end{bmatrix}.\label{eq:full_ffsm_dynamics}\end{equation}
\subsection{Momentum Conservation and the Generalized Jacobian}
\noindent\textbf{Assumption 1 (zero initial momentum):} The initial total linear and angular momenta are zero.
\begin{equation}H_b(q)\dot x_b+H_{bm}(q)\dot q=0.\label{eq:momentum_constraint}\end{equation}
\noindent\textbf{Assumption 2 (invertible base inertia):} $H_b(q)$ is uniformly positive definite in the workspace. Hence
\begin{equation}\dot x_b=-H_b^{-1}(q)H_{bm}(q)\dot q\triangleq J_{bm}(q)\dot q.\label{eq:base_joint_jacobian}\end{equation}
The end-effector velocity satisfies
\begin{equation}\dot x_e=J_b(q)\dot x_b+J_m(q)\dot q.\label{eq:end_velocity_common}\end{equation}
Thus
\begin{equation}\dot x_e=[J_m(q)+J_b(q)J_{bm}(q)]\dot q\triangleq J_g(q)\dot q.\label{eq:generalized_jacobian}\end{equation}
Near rank deficiency, velocity amplification, stronger base reaction, and reduced manipulability may occur simultaneously.
\subsection{Equivalent Joint-Space Dynamics}
Differentiating the momentum constraint and eliminating $\ddot x_b$ gives
\begin{equation}M_e(q)\ddot q+C_e(q,\dot q)\dot q=\tau+d(t),\label{eq:joint_model}\end{equation}
where
\begin{equation}M_e(q)=H_m(q)-H_{bm}^\top(q)H_b^{-1}(q)H_{bm}(q).\label{eq:generalized_inertia}\end{equation}
Here $C_e\dot q$ collects velocity-dependent terms and $d(t)$ denotes lumped parameter errors, unmodeled dynamics, external disturbances, and actuator errors. Assume $M_e$ is uniformly positive definite. Relabeling the joint variable as $q_m$ gives the basic model
\begin{equation}M_e(q_m)\ddot q_m+C_e(q_m,\dot q_m)\dot q_m=\tau+d(t).\label{eq:joint_model_relabelled}\end{equation}
This reduction retains the equivalent effect of base reaction and provides the basis for computed-torque compensation.
\begin{equation}H_b\ddot x_b+H_{bm}\ddot q+\dot H_b\dot x_b+\dot H_{bm}\dot q=0.\label{eq:momentum_derivative}\end{equation}
The equivalent inertia is the Schur complement and is positive definite under the stated assumptions.

\begin{remark}[Numerical spatial CRBA implementation]
The full inertia matrix can be computed using the spatial composite rigid-body algorithm, $H_{\mathrm{full}}=\sum_{i=0}^nJ_i^\top I_iJ_i$, with $O(n^2)$ complexity. The Schur complement and the equivalent Coriolis vector can then be extracted for controller design. The simulations use a positive-definite nominal joint-space model; a complete CRBA recursive implementation is left for future validation.
\end{remark}

\section{Problem Statement}\label{sec:problem}
Given a reference $q_d\in C^2([0,\infty);\mathbb R^n)$ with bounded
$q_d$, $\dot q_d$, and $\ddot q_d$, define
\begin{equation}
 e=q_m-q_d,\qquad \dot e=\dot q_m-\dot q_d,\qquad
 z=[e^\top,\dot e^\top]^\top.
 \label{eq:tracking_errors}
\end{equation}
The objective is to prescribe the total settling time through two design
parameters, $T_o$ and $h$:
\begin{equation}
 T=T_{\rm pre}=T_o+2h.
 \label{eq:prescribed_time_definition}
\end{equation}
\begin{problem}[Prescribed-time joint trajectory tracking]\label{prob:tracking}
Design $\tau(t)$ such that $e(t)=\dot e(t)=0$ for all $t\geq T$ under
ideal continuous-time exact compensation. Under nonidealities, establish an
explicit ultimate bound in terms of the residual acceleration disturbance.
\end{problem}

\subsection{Basic Assumptions}
\begin{assumption}[Equivalent inertia]\label{as:inertia}
There is $\underline m>0$ such that
$M_e(q_m)\succeq\underline m I_n$ on the considered workspace. The nominal
$M_e$ and $C_e$ used by the controller are available.
\end{assumption}
\begin{assumption}[Reference and measurements]\label{as:reference}
$q_d\in C^2$, and $q_d$, $\dot q_d$, and $\ddot q_d$ are bounded. The states
$q_m$ and $\dot q_m$ are measured and all signals used by the controller are
well defined on the considered interval.
\end{assumption}
\begin{assumption}[Delay history]\label{as:delay_history}
The function $z(t)$ is available on $[T_o-h,T_o]$ and the delayed signal is
implemented from a bounded, piecewise-continuous buffer. This history need not
satisfy the nominal model before $T_o$.
\end{assumption}
\begin{assumption}[Disturbance regularity]\label{as:disturbance}
The lumped disturbance is such that $\Delta_a=M_e^{-1}d$ is absolutely
continuous and there exist finite constants $\bar\Delta_a,L_d$ satisfying
\begin{equation}
 \|\Delta_a(t)\|\leq\bar\Delta_a,\qquad
 \|\dot\Delta_a(t)\|\leq L_d
\end{equation}
for almost every $t$. The observer gains are selected using the hypotheses of
the observer theorem stated in Section~\ref{sec:controller}.
\end{assumption}
\begin{assumption}[Ideal implementation]\label{as:ideal}
The exact prescribed-time result refers to continuous-time implementation,
absence of actuator saturation and measurement noise, exact model matching,
and an observer satisfying Assumption~\ref{as:observer_interface}. With
sampling, saturation, or model error, only the practical bound in
Section~\ref{sec:practical} is claimed.
\end{assumption}

\subsection{Delay-system preliminaries}
For a retarded linear delay system
\begin{equation}
 \dot x(t)=F(t)x(t)+G(t)x(t-h),
 \label{eq:abstract_delay_system}
\end{equation}
with a bounded initial history on an interval of length $h$, let $\Psi(t,s)$
denote its causal state-transition operator. If $F$ and $G$ are $2h$-periodic,
then the sampled state at period boundaries has a linear map
\begin{equation}
 x(2(k+1)h)=\Delta(h)x(2kh),\qquad k=0,1,\ldots,
 \label{eq:period_map_abstract}
\end{equation}
whenever the history has been included in the state representation. For the
PDF design below, $G(t)=0$ on the first half-period, so the delayed signal on
the second half-period is explicitly determined by the first-half trajectory
and the map can be calculated by variation of constants.

\begin{definition}[Single-period nullification]\label{def:single_period_nullification}
A PDF system has single-period nullification if its one-period map satisfies
$\Delta(h)=0$. In that case the state at the end of every complete period is
zero under ideal homogeneous dynamics, independently of the state entering the
period.
\end{definition}
\begin{definition}[Nilpotent nullification]\label{def:nilpotent_nullification}
If $\Delta(h)$ is nilpotent of index $\nu_0$, then the sampled state is zero
at or before $2\nu_0h$ under ideal homogeneous dynamics.
\end{definition}

\section{Controller Design}\label{sec:controller}

This section separates the observer result from the delayed-feedback result. The
latter is proved exactly; the observer is required to satisfy the explicit
finite-time interface stated in Lemma~\ref{lem:observer_interface}.

\subsection{Disturbance-observer interface}\label{sec:ptdo}
Let $\chi=\dot q_m$ and define
\begin{equation}
 \dot\chi=u_o+\Delta_a,\qquad
 u_o=M_e^{-1}(q_m)[\tau-C_e(q_m,\dot q_m)\dot q_m],\qquad
 \Delta_a=M_e^{-1}(q_m)d.
 \label{eq:observer_channel}
\end{equation}
The following assumption is the precise interface needed by the controller.
\begin{assumption}[Observer interface]\label{as:observer_interface}
There are $T_o>0$ and an estimate $\hat\Delta_a$ such that, under the
continuous-time observer, $\hat\Delta_a(t)=\Delta_a(t)$ for all $t\geq T_o$.
Before $T_o$, the observation error is bounded. The estimate is formed as
$\hat d=M_e\hat\Delta_a$.
\end{assumption}

For completeness, the observer used in the simulations is the prescribed-time
homogeneous observer
\begin{equation}
\begin{aligned}
 \dot z_1&=z_2+u_o+\ell_1(t)\,\phi_1(\varepsilon_1/\sigma),\\
 \dot z_2&=\ell_2(t)\,\phi_2(\varepsilon_1/\sigma),
 \qquad \varepsilon_1=\chi-z_1,
\end{aligned}\label{eq:ptdo}
\end{equation}
where
\begin{align}
 \phi_1(x)&=\operatorname{sig}^{1-\eta/2}(x)+
             \operatorname{sig}^{1+\eta/2}(x),\\
 \phi_2(x)&=\operatorname{sig}^{2-\eta}(x)+
             \operatorname{sig}^{2+\eta}(x)+\sigma\operatorname{sgn}(x),
 \end{align}
$0<\eta<1$, and $\ell_1,\ell_2$ are selected according to the
prescribed-time homogeneous-observer design in~\cite{jiangPrescribedtimeDisturbanceObserver2024}.
The gains must be selected using the disturbance regularity bounds required by
that design; in particular, Assumption~\ref{as:disturbance} alone is not
sufficient to establish exact convergence of~\eqref{eq:ptdo}.

\begin{lemma}[Observer interface lemma]\label{lem:observer_interface}
Suppose the hypotheses of the prescribed-time homogeneous-observer theorem
used to select $\ell_1,\ell_2$ hold, including the required bound on
$\dot\Delta_a$. Then the Filippov solution of~\eqref{eq:ptdo} satisfies
$z_2(t)=\Delta_a(t)$ for all $t\geq T_o$.
\end{lemma}
\begin{proof}
The error variables $\varepsilon_1=\chi-z_1$ and
$\varepsilon_2=\Delta_a-z_2$ satisfy the homogeneous observer error system
obtained by subtracting~\eqref{eq:ptdo} from~\eqref{eq:observer_channel}.
The cited prescribed-time homogeneous-observer theorem applies to this
system because its vector field is homogeneous in the two error variables and
its perturbation is $\dot\Delta_a$. It provides a Lyapunov function $V$ and
positive constants $c_1,c_2$ for which the time-rescaled error system obeys
$\dot V\leq-c_1V^{\alpha}-c_2V^{\beta}$, with $0<\alpha<1<\beta$.
The standard fixed-time integral bound gives a settling time no larger than
$T_o$ after the prescribed time transformation. Strong invariance of the
sliding set then gives $\varepsilon_1=\varepsilon_2=0$ for all subsequent
time. Hence $z_2=\Delta_a$ and $\hat d=d$ after $T_o$.
\end{proof}

\subsection{Dynamic compensation and error linearization}
Use
\begin{equation}
 \tau=M_e(q_m)v+C_e(q_m,\dot q_m)\dot q_m-\hat d,
 \qquad v=\ddot q_d+\nu.
 \label{eq:computed_torque}
\end{equation}
Under exact model matching and Assumption~\ref{as:observer_interface},
$\hat d=d$ for $t\geq T_o$, and~\eqref{eq:joint_model_relabelled} gives
$\ddot q_m=v$. Consequently, with
$e=q_m-q_d$ and $z=[e^\top,\dot e^\top]^\top$,
\begin{equation}
 \dot z=Az+B\nu,
 \qquad A=\begin{bmatrix}0&I_n\\0&0\end{bmatrix},
 \qquad B=\begin{bmatrix}0\\I_n\end{bmatrix}.
 \label{eq:linear_error}
\end{equation}
The pair $(A,B)$ is controllable since
$\operatorname{rank}[B,AB]=2n$.

\subsection{Smooth periodic delayed feedback}\label{sec:pdf_design}
Choose $K_0$ so that $A_c=A+BK_0$ is Hurwitz. Let $h>0$ and let
$R_h:[0,2h]\to\mathbb R$ be smooth, nonnegative, zero on $[0,h]$, and
positive on a subset of positive measure in $(h,2h)$. Define
\begin{equation}
 W_c=\int_h^{2h}e^{-A_cs}B R_h(s)B^\top e^{-A_c^\top s}\,ds.
 \label{eq:wc_matrix}
\end{equation}
\begin{lemma}[Positivity of the PDF Gramian]\label{lem:gramian}
If $(A_c,B)$ is controllable, then $W_c\succ0$.
\end{lemma}
\begin{proof}
For any $x\in\mathbb R^{2n}$,
\[
 x^\top W_cx=\int_h^{2h}R_h(s)\|B^\top e^{-A_c^\top s}x\|^2ds\geq0.
\]
If this quantity is zero, continuity implies
$B^\top e^{-A_c^\top s}x=0$ on an interval contained in the set where
$R_h>0$. The left-hand side is analytic in $s$, hence it vanishes for all
$s$. Differentiating at any point gives
$B^\top(A_c^\top)^kx=0$ for every $k\geq0$. Controllability of
$(A_c,B)$ implies $x=0$. Therefore $W_c$ is positive definite.
\end{proof}

For local PDF time $\vartheta=t-T_o$, define the $2h$-periodic gain
\begin{equation}
 K_c(\vartheta)=\begin{cases}
 0,&0\leq\vartheta<h,\\
 R_h(\vartheta)B^\top e^{-A_c^\top\vartheta}W_c^{-1}
 e^{A_c(h-\vartheta)},&h\leq\vartheta<2h,
 \end{cases}
 \label{eq:kc_pdf}
\end{equation}
and extend it periodically. The control input is
\begin{equation}
 \nu(t)=K_0z(t)-K_c(t-T_o)z(t-h),\qquad t\geq T_o.
 \label{eq:pdf_error_input}
\end{equation}
Thus the complete controller is
\begin{equation}
 \tau=M_e\left[\ddot q_d+K_0z-K_c(t-T_o)z(t-h)\right]
 +C_e\dot q_m-\hat d.
 \label{eq:final_controller}
\end{equation}

\section{Stability and Convergence Analysis}\label{sec:stability}
\subsection{Exact prescribed-time tracking under ideal conditions}
By Theorem~\ref{thm:ptdo}, $\hat d=d$ for $t\ge T_o$ under the stated ideal assumptions. During Phase I, $K_c=0$ and
\begin{equation}\dot z=A_cz+B\varepsilon_2,\qquad \|z(t)\|\le\kappa e^{-\lambda t}\|z(0)\|+\frac{\kappa}{\lambda}\sup_{0\le s\le t}\|\varepsilon_2(s)\|.\label{eq:phaseI_ISS}\end{equation}
Thus the state entering Phase II is bounded. In Phase II, exact compensation gives the PDF delay system of Section~\ref{sec:controller}; its one-period map is zero, so
\begin{proposition}[Phase-I input-to-state stability]\label{prop:phaseI_ISS}
The preceding bound holds and $z(t)$ is bounded on $[0,T_o]$.
\end{proposition}
\begin{proof}
Variation of constants for the Hurwitz matrix $A_c$ gives the displayed inequality.
\end{proof}
\begin{equation}e(t)=0,\quad\dot e(t)=0,\qquad t\ge T_o+2h.\label{eq:main_ideal_result}\end{equation}
Therefore $T_{\mathrm{pre}}=T_o+2h$, independently of the initial error. If the map is nilpotent of index $\nu_0$, the bound becomes $T_o+2\nu_0h$.
\begin{theorem}[Exact prescribed-time tracking]\label{thm:main_ideal}
Under the ideal conditions, $e(t)=\dot e(t)=0$ for all $t\ge T_o+2h$.
\end{theorem}
\begin{proof}
Phase I is bounded by the proposition. After $T_o$, exact compensation yields the PDF system, and Theorem~\ref{thm:pdf_gain} gives $z=0$ after one period.
\end{proof}
\begin{corollary}[Prescribed-time upper bound]\label{cor:prescribed_time}
The convergence bound is $T_o+2h$ and is independent of the initial error.
\end{corollary}

\subsection{Practical Bounded Tracking}\label{sec:practical}
Define the residual acceleration disturbance
\begin{equation}\Delta_r=M_e^{-1}[\Delta d_o+\Delta_m+\Delta_\tau].\label{eq:residual_acceleration}\end{equation}
If $\|\Delta_r\|\le\bar\Delta_r$, then
\begin{equation}\dot z=A_cz-BK_c(t)z(t-h)+B\Delta_r(t).\label{eq:perturbed_delay_error}\end{equation}
Let $\Gamma_h$ denote the maximum one-period input-to-state integral gain.
\begin{assumption}[Bounded residual disturbance]\label{as:residual_bounded}
There is $\bar\Delta_r\ge0$ such that $\|\Delta_r(t)\|\le\bar\Delta_r$ for all $t\ge T_o$.
\end{assumption}
\begin{definition}[PDF periodic disturbance gain]\label{def:pdf_gain}
$\Gamma_h$ is the maximum one-period input-to-state integral gain of the homogeneous delayed system.
\end{definition} At $t_k=T_o+2kh$, exact PDF nullification gives
\begin{theorem}[Practical bounded tracking at period boundaries]\label{thm:practical_boundary}
\begin{equation}\|z(t_{k+1})\|\le\Gamma_h\bar\Delta_r.\label{eq:pdf_boundary_bound}\end{equation}
\end{theorem}
\begin{proof}
The homogeneous one-period term vanishes. Variation of constants and the residual bound give the inequality.
\end{proof}
\begin{corollary}[Bounded tracking at arbitrary times]\label{cor:arbitrary_time}
Under the preceding assumptions, there is a constant $\Gamma_h^*>0$ such that $\|z(t)\|\le\Gamma_h^*\bar\Delta_r$ for all $t\ge T_o$.
\end{corollary}

If instead $\|\Delta_K(h)\|\le\varepsilon_\Delta<1$, then
\begin{equation}\|z_{k+1}\|\le\varepsilon_\Delta\|z_k\|+\Gamma_h\bar\Delta_r,\qquad \limsup_{k\to\infty}\|z_k\|\le\frac{\Gamma_h\bar\Delta_r}{1-\varepsilon_\Delta}.\end{equation}
If $\|M_e^{-1}\|\le\bar m^{-1}$ and the three residual terms are bounded by $\bar d_o,\bar\Delta_m,\bar\Delta_\tau$, then
\begin{equation}\bar\Delta_r\le\bar m^{-1}(\bar d_o+\bar\Delta_m+\bar\Delta_\tau).\label{eq:explicit_delta_bar}\end{equation}
For task-space tracking, a damped generalized-Jacobian inverse generates the joint reference:
\begin{equation}\dot q_d=J_g^\#[\dot x_e^d-K_x(x_e-x_e^d)],\quad J_g^\#=J_g^\top(J_gJ_g^\top+\lambda^2I)^{-1}.\label{eq:task_space_velocity}\end{equation}
Gramian ill-conditioning at high dimension can make the ideal nullification approximate; regularization or truncated SVD is recommended.
\begin{proposition}[Approximate nullification bound]\label{prop:approx_null}
If $\|\Delta_K(h)\|\le\varepsilon_\Delta<1$, then
\begin{equation}\|z_{k+1}\|\le\varepsilon_\Delta\|z_k\|+\Gamma_h\bar\Delta_r.\end{equation}
Thus $\limsup_{k\to\infty}\|z_k\|\le\Gamma_h\bar\Delta_r/(1-\varepsilon_\Delta)$.
\end{proposition}
\begin{corollary}[Explicit error bound]\label{cor:explicit_bound}
If the residual terms have bounds $\bar d_o,\bar\Delta_m,\bar\Delta_\tau$, then the bound is obtained from
\begin{equation}\bar\Delta_r\le\bar m^{-1}(\bar d_o+\bar\Delta_m+\bar\Delta_\tau).\end{equation}
\end{corollary}

\section{Simulation Verification}\label{sec:simulation}
A seven-degree-of-freedom nominal positive-definite joint-space model is used to evaluate numerical feasibility. The step size is $T_s=10^{-3}\,\mathrm{s}$ and the duration is $10\,\mathrm{s}$. A fifth-order polynomial reference is generated over $8\,\mathrm{s}$. The delay is $h=1\,\mathrm{s}$, with $R_h(\theta)=\sin^4(\pi(\theta-h)/h)$ on the second half-period. The gains are $K_p=\operatorname{diag}(25,25,25,22,22,20,20)$ and $K_d=\operatorname{diag}(10,10,10,9,9,8,8)$, with $K_0=[-K_p\ -K_d]$. Disturbances are $d_i=0.15\sin(0.7t+0.3i)+0.05\cos(1.3t+0.2i)$. Observer parameters are $T_o=T_c=1\,\mathrm{s}$, $\eta=0.3$, and $\sigma=0.4$.

\subsection{Results and Discussion}
The simulations show bounded base reaction motion, end-effector motion, joint tracking, and control torques under disturbance. The Gramian condition number is $5.796\times10^5$; consequently PDF and PD have nearly identical performance for this 14-state system. Without disturbance, terminal errors are about $10^{-9}$ rad. With disturbance, PD/PDF terminal error is approximately $1.4\times10^{-4}$ rad, whereas PDF+PTDO gives $3.327\times10^{-4}$ rad and a peak torque of $1920.6$ N m, indicating observer overcompensation under the selected discrete-time parameters. The result motivates larger $T_c$, parameter tuning, and saturation limits.

\begin{figure*}[t]
\centering
\subfloat[Free-floating trajectory\label{fig:free_floating_trajectory}]{\includegraphics[width=0.48\textwidth]{figures/free_floating_trajectory.pdf}}
\hfill
\subfloat[Base position\label{fig:base_position}]{\includegraphics[width=0.48\textwidth]{figures/base_position.pdf}}\\[1ex]
\subfloat[End-effector position\label{fig:end_effector_position}]{\includegraphics[width=0.48\textwidth]{figures/end_effector_position.pdf}}
\hfill
\subfloat[Joint posture\label{fig:joint_posture}]{\includegraphics[width=0.48\textwidth]{figures/joint_posture.pdf}}
\caption{Free-floating base, end-effector, and joint trajectories. Each panel is exported independently by MATLAB and assembled using LaTeX.}
\label{fig:motion}
\end{figure*}

\begin{figure*}[t]
\centering
\subfloat[Joint position tracking\label{fig:joint_position}]{%
\includegraphics[width=0.48\textwidth]{figures/joint_position_tracking.pdf}}
\hfill
\subfloat[Joint velocity tracking\label{fig:joint_velocity}]{%
\includegraphics[width=0.48\textwidth]{figures/joint_velocity_tracking.pdf}}\\[1ex]
\subfloat[Position and velocity tracking errors\label{fig:tracking_errors}]{%
\includegraphics[width=0.48\textwidth]{figures/tracking_errors.pdf}}
\hfill
\subfloat[Control torque\label{fig:control_torque}]{%
\includegraphics[width=0.48\textwidth]{figures/control_torque.pdf}}
\caption{Joint tracking results: desired versus actual joint position and velocity, tracking errors, and control torque.}
\label{fig:joint_summary}
\end{figure*}

\begin{table}[t]\centering\caption{Simulation metrics}\label{tab:sim_metrics}\begin{tabular}{cccc}\toprule Controller and condition & $\|e(T)\|$ [rad] & $\mathrm{RMS}(e)$ [rad] & $\max|\tau_i|$ [N m]\\\midrule PD, no disturbance & $1.333\times10^{-9}$ & $5.039\times10^{-10}$ & 258.9\\ PDF, no disturbance & $1.333\times10^{-9}$ & $5.039\times10^{-10}$ & 258.9\\ PD, disturbance & $1.373\times10^{-4}$ & $5.188\times10^{-5}$ & 258.9\\ PDF, disturbance & $1.373\times10^{-4}$ & $5.188\times10^{-5}$ & 258.9\\ PDF+PTDO, disturbance & $3.327\times10^{-4}$ & $1.258\times10^{-4}$ & 1920.6\\\bottomrule\end{tabular}\end{table}

\begin{figure*}[t]
\centering
\subfloat[True and estimated disturbance\label{fig:observer_estimate}]{\includegraphics[width=0.48\textwidth]{figures/observer_estimates.pdf}}
\hfill
\subfloat[Observer error\label{fig:observer_error}]{\includegraphics[width=0.48\textwidth]{figures/observer_error.pdf}}
\caption{Disturbance-observer response.}
\label{fig:observer}
\end{figure*}
\section{Conclusions and Future Work}\label{sec:conclusion}
This paper developed a prescribed-time controller based on smooth periodic delayed feedback for joint-space tracking of a pre-capture free-floating space manipulator. The four-stage structure combines disturbance observation, dynamic compensation, error linearization, and periodic delayed feedback. Under continuous-time implementation, an exact model, no saturation, and exact compensation, tracking errors vanish after $T_{\mathrm{pre}}=T_o+2h$. With residual disturbances and approximate nullification, explicit practical bounds were established. A Gramian-based gain construction was presented, and its sensitivity in high-dimensional systems was demonstrated numerically.

Future work includes: (1) a complete spatial-CRBA floating-base simulation; (2) low-sensitivity Gramian shaping and structured delay-feedback design; (3) integration of task-space planning, singularity avoidance, joint limits, and dynamic constraints; (4) discrete-time boundedness theory under sampling, noise, and saturation; and (5) hardware-in-the-loop or air-bearing-table experiments.



\bibliographystyle{IEEEtran}
\bibliography{refs/local}
\end{document}
